\chapter{Division and Multiples}\label{ch:g5:division}

Long division learned in \cref{ch:g4:division} gets two upgrades:
two-digit divisors, and quotients that continue past the point ---
$3 \div 4$ finally gets an answer, $0.75$. The chapter ends with
multiples, divisors and the first divisibility shortcuts.

\section{Dividing by a two-digit number}

\begin{method}[Two-digit divisors]\label{met:g5:division:twodigits}
Same method as before, with one extra skill: guessing how many times
the divisor fits. Write a small table of multiples of the divisor
first. For $986 \div 23$ (multiples of 23: 23, 46, 69, 92, 115,
138, 161, 184, 207):
\begin{enumerate}
  \item $98$ tens: $23 \times 4 = 92$ fits, $23 \times 5 = 115$ does
        not: digit $4$, remainder $6$;
  \item bring down the $6$: $66$: $23 \times 2 = 46$ fits: digit
        $2$, remainder $20$;
  \item quotient $42$, remainder $20$. Check:
        $23 \times 42 + 20 = 966 + 20 = 986$.
\end{enumerate}
\end{method}

\section{Decimal quotients}

\begin{example}[The division that refuses to stop at the remainder]\label{ex:g5:division:decimal}
Share $3$ pizzas among $4$ children: $3 \div 4$ has quotient $0$ and
remainder $3$ --- useless! Continue the division past the point:
\begin{enumerate}
  \item $3$ units $= 30$ tenths; $30 \div 4$: $7$ tenths, remainder
        $2$ tenths;
  \item $2$ tenths $= 20$ hundredths; $20 \div 4 = 5$ hundredths,
        remainder $0$.
\end{enumerate}
So $3 \div 4 = 0.75$: each child gets $0.75$ pizza --- three
quarters, as the fraction language already knew
(\cref{prop:g5:fractions:decimal}).
\end{example}

\begin{method}[Continuing a division past the point]\label{met:g5:division:continue}
When the units are exhausted and a remainder is left:
\begin{enumerate}
  \item write the decimal point in the quotient;
  \item append a zero to the remainder (units become tenths, tenths
        become hundredths\,\dots) and keep dividing;
  \item stop when the remainder is $0$ --- or when the digits start
        repeating forever, like $1 \div 3 = 0.333\dots$: then give a
        rounded value.
\end{enumerate}
\end{method}

\begin{example}\label{ex:g5:division:continue}
$27 \div 6$: quotient $4$, remainder $3$; point; $30$ tenths
$\div\ 6 = 5$: so $27 \div 6 = 4.5$. \quad
$22 \div 8$: $2$, remainder $6$; then $60 \div 8 = 7$ r $4$;
then $40 \div 8 = 5$: so $22 \div 8 = 2.75$.
\end{example}

\section{Multiples and divisors}

\begin{definition}[Multiple, divisor]\label{def:g5:division:multiple}
The \emph{multiples}\index{multiple} of $6$ are its table continued
forever: $0, 6, 12, 18, 24, \dots$ When $24$ is a multiple of $6$,
one also says that $6$ is a \emph{divisor}\index{divisor} of $24$,
or that $24$ is \emph{divisible} by $6$: the division
$24 \div 6$ leaves no remainder.
\end{definition}

\begin{proposition}[Divisibility shortcuts]\label{prop:g5:division:criteria}
Without dividing, a number is divisible:
\begin{itemize}
  \item by $2$ when its last digit is $0$, $2$, $4$, $6$ or $8$
        (the even numbers);
  \item by $5$ when its last digit is $0$ or $5$;
  \item by $10$ when its last digit is $0$;
  \item by $3$ when the \emph{sum of its digits} is divisible by
        $3$ (e.g.\ $741$: $7 + 4 + 1 = 12$, divisible --- and indeed
        $741 = 3 \times 247$).
\end{itemize}
\end{proposition}
\admitted

\begin{omfigure}
\begin{tikzpicture}[scale=0.5]
  \foreach \n in {0,...,29} {
    \pgfmathtruncatemacro{\r}{int(\n/10)}
    \pgfmathtruncatemacro{\c}{mod(\n,10)}
    \pgfmathtruncatemacro{\lbl}{\n+1}
    \pgfmathtruncatemacro{\mtwo}{mod(\lbl,2)}
    \pgfmathtruncatemacro{\mthree}{mod(\lbl,3)}
    \ifnum\mtwo=0
      \fill[omDef!25] (\c,-\r) ++(-0.45,-0.45) rectangle ++(0.9,0.9);
    \fi
    \ifnum\mthree=0
      \draw[omThm, thick] (\c,-\r) circle (0.42);
    \fi
    \node[font=\footnotesize] at (\c,-\r) {\lbl};
  }
  \node[font=\small, below] at (4.5,-2.8)
    {blue squares: multiples of $2$;\ \ red circles: multiples of
    $3$;\ \ both: multiples of $6$};
\end{tikzpicture}

{\small The numbers $1$ to $30$: multiples of $2$ and of $3$ make
patterns --- and the numbers wearing both marks ($6$, $12$, $18$,
$24$, $30$) are exactly the multiples of $6$.}
\end{omfigure}

\begin{example}\label{ex:g5:division:criteria}
Is $534$ divisible by $2$? Ends in $4$: yes. By $5$? No. By $3$?
$5 + 3 + 4 = 12$: yes. So $534$ is divisible by $6$ as well (by $2$
\emph{and} by $3$) --- check: $534 = 6 \times 89$.
\end{example}

\section{Exercises}

\begin{exercise}[$\star$]\label{exo:g5:division:1}
Compute the long divisions (table of multiples first):
$851 \div 23$; \quad $704 \div 32$.
\end{exercise}

\begin{exercise}[$\star$]\label{exo:g5:division:2}
Compute the exact decimal quotients: $9 \div 2$; \quad $27 \div 4$;
\quad $33 \div 5$; \quad $21 \div 8$.
\end{exercise}

\begin{exercise}[$\star$]\label{exo:g5:division:3}
Four friends share a $54$ restaurant bill equally. How much does
each pay? (Continue past the point.)
\end{exercise}

\begin{exercise}[$\star$]\label{exo:g5:division:4}
Compute $10 \div 3$, continuing three digits past the point. What do
you observe? Give the value rounded to the hundredth.
\end{exercise}

\begin{exercise}[$\star$]\label{exo:g5:division:5}
List: the multiples of $7$ up to $70$; the divisors of $18$; the
divisors of $24$.
\end{exercise}

\begin{exercise}[$\star$]\label{exo:g5:division:6}
Divisible by $2$? by $5$? by $10$? by $3$? Test each shortcut on:
$470$; \quad $735$; \quad $8\,001$; \quad $1\,314$.
\end{exercise}

\begin{exercise}[$\star$]\label{exo:g5:division:7}
Find all the numbers between $60$ and $80$ that are divisible by
$3$ \emph{and} by $5$. (Which single divisibility does that
combine?)
\end{exercise}

\begin{exercise}[$\star$]\label{exo:g5:division:8}
A ribbon of $7.2$ m is cut into $6$ equal pieces. How long is each
piece? (Divide a decimal by a whole: share the meters, then the
tenths.)
\end{exercise}

\begin{exercise}[$\star$]\label{exo:g5:division:9}
$1\,000$ marbles are packed in bags of $24$. How many full bags, and
how many marbles remain? (Two-digit divisor.)
\end{exercise}

\begin{exercise}[$\star\star$]\label{exo:g5:division:10}
Can $5$ friends share $7$ chocolate bars equally, with nothing left?
Give each share as a decimal. Same question for $7$ friends sharing
$5$ bars --- what goes differently?
\end{exercise}

\begin{exercise}[$\star\star$]\label{exo:g5:division:11}
Using the digit-sum shortcut, find the smallest digit $d$ that makes
$52d$ (a three-digit number ending in $d$) divisible by $3$. Then
find all such digits $d$.
\end{exercise}
